\documentclass[10pt,twocolumn]{article}
\usepackage[letterpaper,margin=0.75in,columnsep=0.25in]{geometry}
\usepackage{amsmath,amssymb,booktabs,array,microtype,xcolor,hyperref}
\definecolor{ink}{RGB}{22,39,57}
\definecolor{accent}{RGB}{18,105,116}
\definecolor{muted}{RGB}{95,109,120}
\definecolor{pale}{RGB}{237,245,246}
\hypersetup{colorlinks=true,urlcolor=accent,citecolor=accent,linkcolor=ink,
 pdftitle={ScaleRLT: I/O-Aware End-to-End Reinforcement Learning for Recurrent Looped Transformers},pdfauthor={0z5a}}
\setlength{\parindent}{1em}
\setlength{\parskip}{0pt}
\setlength{\textfloatsep}{11pt plus 2pt minus 2pt}
\setlength{\floatsep}{9pt plus 2pt minus 2pt}
\setlength{\tabcolsep}{4pt}
\renewcommand{\arraystretch}{1.08}
\makeatletter
\renewcommand\section{\@startsection{section}{1}{0pt}{-2.3ex plus -0.5ex minus -0.2ex}{1ex plus 0.2ex}{\large\bfseries\color{ink}}}
\renewcommand\subsection{\@startsection{subsection}{2}{0pt}{-1.8ex plus -0.4ex minus -0.2ex}{0.7ex plus 0.2ex}{\normalsize\bfseries\color{ink}}}
\makeatother
\newcommand{\pass}{\textsc{Pass}}
\newcommand{\fail}{\textsc{Fail}}
\newcommand{\pending}{\textsc{Not run}}
\newcommand{\code}[1]{\texttt{#1}}
\title{ScaleRLT: I/O-Aware End-to-End Reinforcement Learning for Recurrent Looped Transformers}
\author{0z5a}
\date{Research draft v0.1, October 6, 2026}
\begin{document}
\twocolumn[{
\begin{center}
{\LARGE\bfseries\color{ink}ScaleRLT: I/O-Aware End-to-End Reinforcement Learning\\[3pt]for Recurrent Looped Transformers\par}
\vspace{9pt}
{\large 0z5a\par}
\vspace{5pt}
{\small\color{muted}Research draft v0.1\quad $\cdot$\quad October 6, 2026\par}
\vspace{12pt}
\begin{minipage}{0.94\textwidth}\small
\textbf{Abstract.}
Looped transformers share physical weights across recurrent invocations, but reinforcement learning (RL) still repeats prefix computation, retains invocation-specific states, and moves scores, policies and optimizer checkpoints. We present ScaleRLT, an I/O-aware execution prototype that preserves declared looped-policy and gradient contracts while coordinating prefix adjoints, suffix batching, rematerialization, stage-specific KV storage and output export. Across three tiny FP32 CPU model families, a prefix-heavy fixed-trace update achieves $2.96$--$3.68\times$ speed ratios with complete gradient and update checks; shorter-prefix and single-response controls expose regressions. Stage-specific Huginn R32 allocation reduces measured tiny KV tensor storage by $48.44\%$. Twelve tiny H20/Megatron configurations pass paired two-step AdamW checks, including all parameter gradients, moments and accounting. A separately audited official Ouro inference baseline gains $1.142\times$ fixed-work throughput by reducing decode depth, while its development accuracy remains statistically uncertain; it is not an exact-schedule RL gain. BF16 and IQuest strict-update failures are retained. These results establish component feasibility and numerical boundaries. Complete official-model online RL, fresh-worker recovery, high-load stability and independent-seed reward convergence remain open. The principal endpoint is time and allocated GPU-hours to held-out reward quality.
\end{minipage}
\vspace{10pt}\par
{\color{accent}\rule{\textwidth}{0.8pt}}\par
\end{center}
\vspace{4pt}
}]

\section{Introduction}
Recurrent looped transformers offer a distinct way to scale reasoning: apply shared physical parameters repeatedly in latent space rather than store a separate layer stack for every computational step~\cite{ouro,huginn}. This economy of weights does not make their execution states interchangeable. Each invocation has its own hidden state, attention history and gradient contribution. In RL, multiple sampled responses to one prompt add repeated prefix work, reference evaluation, policy publication and persistent optimizer state to that recurrent graph.

A useful optimization must preserve the policy being trained and the evidence that it learns. Faster token generation can change the decode depth or sampling distribution. A correct forward pass can still yield a different optimizer update. An incremented policy version can coexist with zero advantages and no reward-driven learning. We therefore organize ScaleRLT around an accuracy-first question: which execution costs can be reduced while preserving a declared looped objective, and do those changes lower the cost of reaching a fixed held-out reward?

Our prototype extends the existing VIME/Megatron learner and native vllm-rlt backend~\cite{vime,rlt}. Physical parameter sharing, model adapters, recurrent trace identities and publication machinery are inherited infrastructure. Ordinary three-phase differentiable prefix reuse is also prior work~\cite{prefix}. The present engineering focus is their interaction with repeated loops, first-response readout, model stages, actual latent identity and optimizer scheduling.

This initial manuscript contributes an explicit looped-RL execution contract, an implementation spanning prefix and state lifetimes, and an audited evaluation that distinguishes exact transformations from approximate inference baselines. Its evidence includes both positive component results and retained negative controls. It does not yet establish a production-scale scheduling advantage over generic prefix reuse, a complete RL speedup, or reward convergence. Those conclusions require the frozen comparisons described in Section~\ref{sec:quality}.

\section{Looped RL Execution Contracts}
\subsection{Objective and logical update}
Let $\theta$ denote unique physical parameters, $r$ a recurrent invocation, and $i,t$ a trajectory and response position. The actor exposes
\begin{equation}
 \ell_{i,t,r}=\log p_{\theta}^{(r)}(y_{i,t}\mid x_i,y_{i,<t}).
\end{equation}
For response masks $m_{i,t}$, fixed group-relative advantages $\hat A_i$, and declared loop weights $w_{i,t,r}$, the policy-gradient primitive is
\begin{equation}\label{eq:objective}
 L_{\rm PG}=-\mathcal R_{i,t}\!\left[m_{i,t}\hat A_i\sum_r w_{i,t,r}\ell_{i,t,r}\right].
\end{equation}
The weighted sum of log probabilities is not the log probability of a mixture. The public RLTT function implements an unclipped policy-gradient surrogate with a masked token mean~\cite{rltt}. A response-mean profile is a separately declared objective. Reproducing the original experiment also requires its surrounding minibatch, accumulation, reference, optimizer and reward contracts; matching Equation~\ref{eq:objective} alone is insufficient.

A logical update fixes its reduction denominator before splitting suffix waves or microbatches. Parameters remain unchanged until all contributions have been accumulated. The ordinary MCore schedule divides microbatch losses by the microbatch count; its adapter compensates for that division. The prefix path consumes the logical objective directly, applies optimizer loss scaling once to suffix backward, and injects already-scaled prefix adjoints without a second factor. Empty and masked responses retain their original accounting identities.

\subsection{Identity and model-specific state}
Every execution node carries the policy digest, checkpoint revision, family, physical layer, loop, token range and execution profile. Sharing a weight tensor does not authorize sharing distinct loop KV. Raw and sampling-processed scores remain separate fields; sampled traces preserve depth, stop reason, sampling transforms, publication generation and latent identity.

Ouro and Nanbeige permit compatible fixed-depth text-prefix groups. Huginn additionally conditions recurrent execution on a recorded latent seed and profile. Equal prompt tokens with independent latents are not equal recurrent prefixes. Compatible latent groups use the shared graph; singletons execute the complete replay. A backward receipt records the actual shared fraction instead of assuming all grouped prompts can reuse recurrent states.

The first response token is predicted from the final prompt hidden state. That readout is a direct cross-boundary gradient path at each supervised loop. Huginn adds prelude injection and coda dependencies. Dropping these edges can preserve many suffix scores while silently changing the update.

\section{I/O-Aware System Design}\label{sec:design}
\begin{figure*}[t]
\centering
\setlength{\unitlength}{1pt}
\begin{picture}(480,112)
\put(0,64){\color{accent}\framebox(104,36){\shortstack{\bfseries Committed policy\\\small Rollout + reward}}}
\put(124,64){\color{accent}\framebox(104,36){\shortstack{\bfseries Shared prefix\\\small All loop boundaries}}}
\put(248,64){\color{accent}\framebox(104,36){\shortstack{\bfseries Suffix waves\\\small Readout + adjoints}}}
\put(372,64){\color{accent}\framebox(104,36){\shortstack{\bfseries Logical update\\\small Adam + publication}}}
\put(106,82){\vector(1,0){16}}
\put(230,82){\vector(1,0){16}}
\put(354,82){\vector(1,0){16}}
\put(300,61){\vector(0,-1){18}}
\put(176,43){\vector(0,1){18}}
\put(176,43){\line(1,0){124}}
\put(192,28){\makebox(112,10){\small Summed boundary gradients}}
\put(0,1){\color{muted}\framebox(476,20){\small Policy/latent identity\quad $\cdot$\quad Stage KV + rematerialization\quad $\cdot$\quad Verified durable state}}
\end{picture}
\caption{ScaleRLT execution boundaries. All suffix waves contribute before one prefix backward and one logical optimizer update. The state layer tracks model-specific dependencies and durable policy/optimizer artifacts. This is a design diagram, not a measured overlap timeline.}
\label{fig:system}
\end{figure*}

\subsection{Differentiable prefix scheduling}
Let $z(\theta)$ contain all shared boundary states and $L_i$ the appropriately weighted suffix loss. Over real arithmetic,
\begin{equation}\label{eq:prefix}
 \nabla_\theta L=\sum_i g_{\theta,\mathrm{suffix}_i}
       +J_{\theta,z}^{\mathsf T}\sum_i g_{z,i}.
\end{equation}
One prefix forward captures the live graph. Suffix waves read detached boundary leaves, execute their private causal histories and accumulate boundary adjoints. A final prefix backward injects the sum into the live graph. Equation~\ref{eq:prefix} follows the established three-phase schedule~\cite{prefix}; ScaleRLT extends its boundary representation to every relevant loop and stage, including the first-response hidden state.

The first correctness reference executes suffix projections individually and is slower than complete packed replay. The subsequent iteration batches projections, MLPs, normalization and readout within each suffix wave while retaining segmented causal attention. Attention visibility, positions, response masks and latent identities remain unchanged. The update is verified against an independent complete replay, rather than only against another prefix implementation.

\subsection{Rematerialization and state lifetime}
The prototype exposes loop checkpoints, physical-layer intervals and token query chunks. Query chunks retain the complete causal key/value prefix and all corresponding gradients. The joint Ouro/Nanbeige path checkpoints the prompt producer and full-depth suffix waves while preserving summed prefix adjoints. Token and label tensors are explicit checkpoint inputs; retained mutation failures showed that capturing mutable response tensors in closures can corrupt backward replay.

The measured CPU plan selector admits only gradient-qualified records for the same workload, timing scope and storage metric. A full GPU residency search, sharded weight windows and arbitrary joint layer/query subdivision remain future implementation. Saved forward storage does not represent backward workspace, allocator peaks or total learner memory. Huginn joint prefix rematerialization remains unsupported.

\subsection{Readout and output traffic}
Response-position readout avoids projecting unused prompt positions. Vocabulary tiling computes selected-token scores by exact log-sum-exp without retaining a full loop--token--vocabulary output. Entropy and categorical KL require their complete distributions and reductions; substituting selected-token or top-$k$ scores would change the objective. The first-order readout primitive is numerically tested, while complete online full-KL and learned-credit integration remain pending.

The inherited offline generator repeatedly materializes cumulative selected-score lists although its consumer needs final outputs. For $T$ output tokens, the cumulative path exports $T(T+1)/2$ score elements; terminal-only export materializes $T$. This is a count of objects, not an end-to-end speed prediction. Streaming preserves its cumulative interface. Tokens, depth, score semantics, stop/cancel behavior and speculative acceptance ordering remain part of the output contract.

\subsection{Stage-specific KV and rollout forks}
For a prelude and coda executed once around a recurrent core,
\begin{equation}\label{eq:kv}
 N_{\rm KV}=L_{\rm pre}+R L_{\rm core}+L_{\rm post}.
\end{equation}
Huginn R32 therefore needs $2+32\cdot4+2=132$ logical KV planes instead of a rectangular $32\cdot8=256$. Separate core and boundary pools allocate, grow and release complete page bundles. Byte budgets round down by the exact integer bundle size.

An independent cold-group rollout fork shares a leader's compatible prompt KV and final prompt hidden state. Followers retain private sampling generators; partial tails become private on continuation writes. Admission counts snapshot ownership and may reclaim snapshots in favor of recomputation. This opt-in mechanism is CPU-qualified for Ouro/Nanbeige. Huginn recurrent-prefix sharing and official CUDA stream-ordering qualification are still open.

\subsection{Policy and checkpoint I/O}
Publication and recovery are first-class phases in the complete RL budget. A committed publication records the physical-weight manifest and worker acknowledgments. Recovery must restore actor, immutable reference, optimizer moments/counters, worker RNG and generation identities before reproducing the next update.

A task-owned archiver streams bounded 512 MiB parts with byte counts and SHA-256 identities. A successful part must match both its server digest and an independent download before its temporary local copy is removed. Failed transfers retain their sources; streaming producers must reach natural completion. Small source-archive transfers and a writable draft destination are verified, but complete native optimizer/publication retention has not yet run.

\section{Experimental Method}
\subsection{Evidence scopes and baselines}
We distinguish four experiments: tiny numerical controls, fixed-trace actor timing, fixed-work generation, and fresh online RL. Tiny configurations use actual model-family math but are not official checkpoints. Generation with reduced decode depth is an approximate-policy comparison, not an exact scheduling ablation. Historical H20 results and a repeated control on a different physical H20 retain separate source and device identities.

The learner baseline is complete native VIME/Megatron replay. Additional comparisons comprise a declared source-matched RLTT baseline, ordinary packing/readout/checkpointing, generic differentiable prefix reuse, and Slime/Megatron framework baselines~\cite{rltt,slime,megatron}. Serving baselines include FlashLoop dense/full-method profiles and the optimized vllm-rlt path~\cite{flashloop,rlt}. FlashLoop and Slime common-domain looped qualification and formal matched comparisons remain open. An unsupported profile is reported as unsupported, never as zero throughput.

The core official matrix is Ouro 1.4B/2.6B, Nanbeige4.2 and Huginn-0125. IQuest 40B is a separate extension gated by complete optimizer correctness and actual full-update capacity. Unique checkpoint tensors define physical model size; loop-expanded depth does not create new physical parameters.

\subsection{Timing, memory and numerical gates}
A speed ratio is baseline elapsed time divided by candidate time, or candidate throughput divided by baseline throughput, with identical work and an explicit phase scope. The clean CPU prefix experiment uses three independent process cohorts; approximate 95\% intervals are Student-$t$ intervals on cohort log ratios. Individual repeated steps do not count as independent processes or training seeds.

Peak GPU allocated/reserved bytes, resident KV tensor bytes, live page counts and forward saved storage are distinct metrics. No conversion between them is assumed. Full RL wall time includes rollout, reward, reference/critic, learner, publication and scheduled checkpoint work on the actual critical path. Allocated GPU cost includes idle allocated devices and counts a co-located learner/rollout GPU once.

Numerical qualification checks all trainable parameters, complete gradients, nonzero parameter deltas, Adam first/second moments, counters and logical sample weights. Original failures and their thresholds remain retained. Full updates are not inferred from small output differences or optimizer-state agreement alone.

\section{Measured Results}\label{sec:results}
\subsection{Prefix-heavy actor updates}
Table~\ref{tab:prefix} compares the packed-prefix iteration with an untouched packed-wave parent on tiny width-8, vocabulary-13 FP32 CPU models. Timing includes zeroing gradients, the RLTT loss, forward/backward and SGD; construction, parameter reset, reference, rollout and publication are excluded. Attention remains segmented. Every timed observation verifies loss, full gradients and a nonzero update. The clean campaign has 216 observations across 18 fresh processes, including a serial-prefix reference arm.

\begin{table}[t]
\centering\small
\caption{Tiny CPU actor update, P256/D8/G32. Times are geometric means; ratios and intervals use three process cohorts. Huginn shares the same recorded latent in this fixed trace.}
\label{tab:prefix}
\begin{tabular}{@{}lrrr@{}}
\toprule
Family & Parent ms & Prefix ms & Ratio [95\% CI]\\
\midrule
Ouro &450.72&152.09&2.963 [2.099, 4.184]\\
Nanbeige &219.31&69.79&3.142 [2.039, 4.842]\\
Huginn &605.58&164.77&3.675 [3.100, 4.358]\\
\bottomrule
\end{tabular}
\end{table}

The improvement is workload-dependent. Table~\ref{tab:controls} retains shorter-prefix and singleton controls. All short P24 intervals include one, and P4/D16/G1 central ratios are below one. Online Huginn reuse cannot be inferred from an artificially shared latent. CPU ratios do not establish official-model GPU acceleration or complete-RL benefit.

\begin{table}[t]
\centering\small
\caption{Clean CPU negative controls for the same implementation. Ratio above one favors packed prefix.}
\label{tab:controls}
\begin{tabular}{@{}llr@{}}
\toprule
Family & P/D/G & Ratio [95\% CI]\\
\midrule
Ouro &4/16/1 &0.444 [0.250, 0.788]\\
Nanbeige &4/16/1 &0.564 [0.470, 0.676]\\
Huginn &4/16/1 &0.494 [0.274, 0.889]\\
Ouro &24/4/32 &0.902 [0.667, 1.218]\\
Nanbeige &24/4/32 &0.841 [0.621, 1.140]\\
Huginn &24/4/32 &0.921 [0.548, 1.549]\\
\bottomrule
\end{tabular}
\end{table}

\subsection{Storage and component regressions}
Measured tiny Huginn R32 KV tensor allocation decreases by $48.44\%$, matching Equation~\ref{eq:kv}. Its CPU generation intervals all include one; this result establishes storage reduction rather than stable speed. Both arms use the same token capacity. Graph pools, official GPU peaks and total learner memory remain unmeasured.

A readout-only CPU ablation retains fewer intermediates but is slower: dense/streamed forward--backward ratios range from 0.414 to 0.722. Frozen rematerialization plans reduce observed saved storage by 78.30--83.43\%, but all three clean timing intervals include one. These results motivate GPU-specific work and prevent interpreting reduced stored bytes as automatic acceleration. Additional complete tables are preserved in Appendix~\ref{app:components}.

\subsection{Actual MCore CUDA optimizer checks}
Twelve tiny FP32 configurations execute real MCore DDP and AdamW on the historical H20: Ouro and Nanbeige ordinary/joint schedules, and Huginn independent/mixed latent schedules, each under token-mean and response-mean reduction. All 24 paired candidate updates pass the original full numerical and accounting gates. Table~\ref{tab:cuda} reports maxima over each four-update family/schedule group.

\begin{table}[t]
\centering\small
\caption{Historical H20 FP32/MCore correctness at source 2ad271d. Errors are absolute maxima; all moments, counters and sample-accounting checks pass. No speed ratio is measured.}
\label{tab:cuda}
\begin{tabular}{@{}lrr@{}}
\toprule
Family / schedule & Gradient error & Parameter error\\
\midrule
Ouro ordinary/joint &$1.1921\!\times\!10^{-7}$&$5.9605\!\times\!10^{-8}$\\
Nanbeige ordinary/joint &$5.9605\!\times\!10^{-8}$&$2.9802\!\times\!10^{-8}$\\
Huginn independent &$7.1526\!\times\!10^{-7}$&$9.3023\!\times\!10^{-8}$\\
Huginn mixed &$1.1772\!\times\!10^{-6}$&$2.0949\!\times\!10^{-7}$\\
\bottomrule
\end{tabular}
\end{table}

The original B-stage source d6229d6 is separately repeated on the epoch2 H20 (GPU 8ee), distinct from the historical EDF device and the newly registered d952 instance. Both reduction configurations pass four paired candidate updates. Maximum gradient/parameter errors are $1.1921\times10^{-7}$/$5.9605\times10^{-8}$; the smallest update norms are 0.027031 and 0.028335. All 20 payloads and their manifest are independently verified after natural completion. The 15.706-second wrapper duration is verification time, not a performance comparison. These controls do not qualify official online reward learning or BF16.

Twelve CPU BF16 configurations remain failed at unchanged gates: all 24 update observations fail parameter-delta checks, and six fail the later gradient check. Maximum gradient-relative-L2 and update errors are $2.74\times10^{-2}$ and $1.77\times10^{-3}$. FP32 qualification does not remove this mixed-precision boundary.

\subsection{Official Ouro inference correctness}
A separately audited LoopCD/e8 baseline uses the pinned ordinary Ouro-1.4B checkpoint. FP32 Torch and Triton each pass all 12 numerical cases, 108 readouts, 96 forced continuations and 21 lifecycle/ownership checks at original absolute/relative bounds of $10^{-4}$. Table~\ref{tab:ouro-numeric} reports complete-scope maxima.

\begin{table}[t]
\centering\small
\caption{Official Ouro-1.4B FP32 inference baseline; no tolerance is relaxed. Selected LP is selected-token log probability.}
\label{tab:ouro-numeric}
\begin{tabular}{@{}lrrr@{}}
\toprule
Backend & Hidden state & Logits & Selected LP\\
\midrule
Torch &$1.907\!\times\!10^{-5}$&$3.528\!\times\!10^{-5}$&$1.907\!\times\!10^{-5}$\\
Triton &$1.907\!\times\!10^{-5}$&$3.433\!\times\!10^{-5}$&$1.335\!\times\!10^{-5}$\\
\bottomrule
\end{tabular}
\end{table}

The first BF16/Triton prefill fails its independently frozen 0.02 absolute/relative hidden-state gate: 998 of 2,048 values mismatch, with maximum error 0.230469. Decode, KV, logit, selected-score and lifecycle checks are not reached. Differing BF16 rounding paths are a static diagnostic hypothesis, not a demonstrated cause.

A separate RTX 5090 terminal-output campaign compares 16 official BF16 Ouro pairs across synchronous, asynchronous single-/multi-stream and core-Graph modes at actual resident concurrency 1/4/16/64. Final tokens, depths, processed scores and metadata agree exactly for 1,360 requests and 10,880 output tokens. The short synthetic prompts and eight-token completions verify output semantics; they are not reasoning accuracy or a long-running stability experiment. No independent timing ratio is reported.

\subsection{Approximate inference: speed and quality}
The official FP32/Triton depth baseline provides a useful accuracy-cost boundary. Its GSM8K development experiment evaluates all six arms on 128 frozen IDs, producing 768 outputs and 82,041 generated tokens under the same strict grader. Table~\ref{tab:gsm} retains every arm and the unchanged one-percentage-point noninferiority decision.

\begin{table}[t]
\centering\small
\caption{Official Ouro GSM8K development accuracy. P4 prefill is fixed. Status is paired against D4/off at the original 1 pp margin; this is inference evaluation, not an RL reward curve.}
\label{tab:gsm}
\begin{tabular}{@{}llrrl@{}}
\toprule
Depth & Guidance & Correct & Accuracy & Status\\
\midrule
D4&Off&80/128&62.50\%&Reference\\
D4&Two-head&83/128&64.84\%&Uncertain\\
D3&Off&79/128&61.72\%&Uncertain\\
D3&Two-head&80/128&62.50\%&Uncertain\\
D2&Off&59/128&46.09\%&Degraded\\
D2&Two-head&62/128&48.44\%&Degraded\\
\bottomrule
\end{tabular}
\end{table}

D3/two-head and D4/off have equal correct totals, but their paired accuracy-difference interval is $[-11.54,+11.54]$ pp. This does not establish noninferiority. D3 remains an exploratory cost candidate; independent 505-math/148-code confirmation is unrun.

A separate fixed-work C32 generation experiment uses P512/T128, a 64 GiB KV budget and actual resident A16/S16. C32 describes the initial/total queue, not 32 simultaneously resident sequences. Each of three arms runs in two fresh processes with five warmups and five measured trials per process: 60 trials, 1,920 requests and 245,760 generated tokens overall.

\begin{table}[t]
\centering\small
\caption{Historical H20 official Ouro fixed-work generation. Ratios are geometric means of two fresh-process mean-throughput ratios. Quality remains uncertain for both D3 arms.}
\label{tab:cost}
\begin{tabular}{@{}lrrr@{}}
\toprule
Arm & Fresh 0 tok/s & Fresh 1 tok/s & Ratio vs D4\\
\midrule
D4/off&143.8722&144.4825&$1.0000\times$\\
D3/off&166.7335&167.5084&$1.1591\times$\\
D3/guided&164.6464&164.7555&$1.1424\times$\\
\bottomrule
\end{tabular}
\end{table}

Guidance is about 1.45\% slower than D3/off. D4/off and D3/off each allocate 74,501,625,856 bytes; D3/guided allocates 128 KiB more, with 2 MiB more reserved. Thus this experiment supplies no VRAM-saving claim. The speed comes from an approximate reduced-depth policy; it is neither exact-schedule actor acceleration nor complete-RL time-to-quality.

\subsection{IQuest strict-update boundary}
The pinned IQuest 40B source and 16 checkpoint headers describe 883 BF16 tensors, 39,794,696,320 unique physical parameters and 79,589,392,640 tensor bytes. Only 100,960 header bytes are read; no official weight tensor body is downloaded. Its two loop passes reuse 80 layers and combine first-loop global KV with second-loop local KV through query gates. Metadata alone is not a 40B execution result.

Tiny original-source FP32 calls expose missing global KV and prefill/decode inconsistencies. A functional replay preserves differentiable first-loop KV, the head-major mixed-attention layout and per-query local window. Across 16 configurations and 32 paired AdamW steps per backend, output and complete-gradient gates pass, but strict full-update gates remain failed (Table~\ref{tab:iquest}). Changing the attention operator family does not qualify this learner.

\begin{table}[t]
\centering\small
\caption{Tiny IQuest replay versus original token-by-token autograd. All backends pass 32 output/gradient checks; the unchanged parameter/update gate remains binding.}
\label{tab:iquest}
\begin{tabular}{@{}lrrr@{}}
\toprule
Backend & Grad. rel. L2 max & Update pass & Fail\\
\midrule
Auto SDPA &$3.219\times10^{-7}$&24/32&8\\
Math SDPA &$3.046\times10^{-7}$&24/32&8\\
Explicit eager &$2.961\times10^{-7}$&26/32&6\\
\bottomrule
\end{tabular}
\end{table}

\section{Reward and Scaled-RL Qualification}\label{sec:quality}
Reward quality is the main endpoint. Before formal runs we freeze the held-out items, grader, sampling budget, noninferiority margin and target quality. Each arm uses at least three independent training seeds. For seed $s$, the first evaluation reaching target $q_\star$ defines
\begin{equation}
 T_{q_\star}^{(s)}=\inf\{t:\hat q_s(t)\ge q_\star\},
\end{equation}
with unreached targets reported as censored, not assigned a successful time. We also report final held-out quality, uncertainty across items and seeds, and allocated device cost
\begin{equation}
 \mathrm{GPUh}=\frac{1}{3600}\sum_{j}\mathrm{allocated\ device\ seconds}_j.
\end{equation}

Training telemetry must expose mixed-reward group frequency, advantage variance, gradient norm, reward-driven parameter delta, optimizer moments/counters, KL, entropy, response length and committed policy identities. A zero-advantage rollout cannot demonstrate learning merely by updating a version number. Evaluation cohorts use one committed policy and fixed prompt/completion seeds; inserting evaluation must not alter training RNG.

The currently frozen official Thinking startup contains three online updates, 96 new training completions and 32 held-out completions, plus three complete checkpoints and four publications. It is prepared but unrun. A three-update startup qualifies lifecycle and nonzero learning signal; it does not establish convergence. Fresh-worker recovery must reproduce the next update with actor, reference, optimizer and RNG restored. Retaining complete native states is required rather than using weights-only recovery.

\begin{table}[t]
\centering\small
\caption{Exit gates for a complete empirical paper. Pending entries carry no numerical estimate.}
\label{tab:remaining}
\begin{tabular}{@{}p{0.47\columnwidth}p{0.45\columnwidth}@{}}
\toprule
Endpoint & Current state\\
\midrule
Tiny three-family CUDA Adam &Pass, historical H20\\
H20 epoch2 B-stage Adam &Pass, separate device\\
Official Thinking online startup &Not run\\
Fresh-worker full-state recovery &Not run\\
Independent-seed reward convergence &Not run\\
DP2/4, TP2/4/8, CP2/4/8 &Not qualified\\
Long-running high-load serving &Not qualified\\
Matched FlashLoop/Slime campaigns &Not run\\
Official GPU peak/capacity frontier &Not measured\\
Complete RL time-to-quality &Not measured\\
\bottomrule
\end{tabular}
\end{table}

\subsection{Standard backend acceptance first}
The immediate prerequisite is the standard VIME \code{train.py}/\code{RolloutManager} path in RFC 465~\cite{acceptance}. Its initial milestone uses pinned Ouro-1.4B, full depth four and ordinary online GRPO without reference, critic or KL. Real rollout/update/publication must advance v0 to v1 to v2, followed by matching fresh-process recovery. Default vLLM import independence, mixed-length trace conversion, failed-publication exclusion, KV version isolation and owned-resource lifecycle remain acceptance requirements. This standard control plane precedes new scheduling features. The prepared Thinking/RLTT startup is a separate follow-up, not a substitute for this gate.

\subsection{Reward-correct scaling}
The first formal scaling figure will use actual 1/2/4/8-GPU allocations crossed with batch/concurrency and sequence length, comparing vanilla VIME/vLLM, fixed-depth RLT and ScaleRLT in their common supported domain. Identical checkpoint, objective, sampling and work budgets are required. For a frozen verifier, the principal goodput is
\begin{equation}
 G_{\rm correct}=\frac{N_{\rm valid,correct}}{T_{\rm complete\ RL}}.
\end{equation}
The numerator counts complete, correctly graded samples with valid committed-policy traces; the denominator includes incorrect samples and all RL phases. Correct-sample token goodput, final reward, time-to-quality, memory and failure tails accompany it. Queued concurrency and allocated devices do not imply resident concurrency or useful work. Current one-device resource metadata supplies no multi-GPU scaling point.

\subsection{Trained loop-native compute policy}
Merged feature collection in PR 454 supplies versioned target features and immutable same-version head snapshots~\cite{draftfeatures}; it does not train a draft or measure speculation acceleration. We will use this chain to train a loop-aware proposal/exit controller after standard-backend acceptance. Proposals verified by a full-depth target and approximate early-exit policies require separate score and quality contracts. Controller training/collection cost and target verification enter the full RL budget. The decisive comparison is higher reward-correct RL goodput under the frozen reward/accuracy noninferiority gate, rather than proposal acceptance rate alone. Controller training, adaptive execution and this comparison are not yet run.

Formal ablations compare the frozen native baseline, prefix-only B, joint-state C, and combined BC under the same objective and natural online sampling. Generic prefix reuse and standard packing/readout/recomputation remain additional controls. Scaled-load experiments preserve total GPU budget and actual resident/queued concurrency, recording goodput, TTFT/TPOT and queue-delay percentiles, failures, staleness, memory and publication latency. Short synthetic output equivalence cannot replace these sustained-load measurements.

\section{Related Work and Limitations}
Ouro and Huginn scale latent computation with shared recurrent weights~\cite{ouro,huginn}; their architectures, depth policies and latent semantics remain model-specific. RLTT supplies a public loop-weighted policy-gradient formulation~\cite{rltt}. ScaleRLT currently implements declared uniform/progressive profiles, not every original private experiment configuration.

Schedule-level shared-prefix reuse already provides prefix forward, suffix microbatches with gradient-cache accumulation, and delayed prefix backward~\cite{prefix}. This is the relevant generic scheduling baseline. A distinct loop-specific benefit would require a measured advantage after ordinary prefix sharing and rematerialization, rather than attributing generic reuse to a new method. FlashLoop studies efficient looped inference via lazy updates~\cite{flashloop}; its approximate-method quality and common supported domain must be measured separately. Megatron and Slime provide the framework context for scalable training and RL~\cite{megatron,slime}.

The current strongest exact-schedule timing evidence is tiny CPU fixed-trace replay, which may be dominated by invocation overhead. Huginn's prefix-heavy case uses artificially shared latent identities. Official CUDA output qualification uses short synthetic prompts and is not a stability campaign. The C32 depth study changes computation and retains uncertain quality. BF16 learner/inference and IQuest complete-update failures prevent wider numerical qualification. No complete original-stack RLTT, Slime or FlashLoop campaign, official-model full update, multi-seed learning curve, or distributed capacity frontier is available in this draft.

\section{Conclusion}
ScaleRLT treats loop identity, prefix gradients, state lifetime and policy/checkpoint I/O as one RL execution problem. The prototype preserves complete FP32 updates in tiny real-MCore CUDA controls and improves prefix-heavy CPU updates, while storage reductions and negative controls expose the costs and limits of each transformation. The next empirical requirement is standard-backend acceptance, reward-correct multi-GPU scaling, and a trained loop-native compute policy with full-state recovery and independent-seed held-out reward. Those measurements will determine whether the system reduces time and GPU-hours to useful quality.

\begin{thebibliography}{10}\small
\bibitem{ouro} R.-J. Zhu et al. Scaling Latent Reasoning via Looped Language Models. \href{https://arxiv.org/abs/2510.25741}{arXiv:2510.25741}, 2025.
\bibitem{huginn} J. Geiping et al. Scaling up Test-Time Compute with Latent Reasoning: A Recurrent Depth Approach. \href{https://arxiv.org/abs/2502.05171}{arXiv:2502.05171}, 2025.
\bibitem{rltt} RLTT contributors. Public loop-weighted policy-gradient implementation. \href{https://github.com/jonwill8/RLTT/blob/06189850bb23a1b1b715ad29768e68f36b1c4a19/rltt_experiments/verl_rltt/rltt_algos.py}{Source revision 0618985}.
\bibitem{prefix} P. Li et al. Schedule-Level Shared-Prefix Reuse for LLM RL Training. \href{https://arxiv.org/html/2606.01143v3}{arXiv:2606.01143v3}, 2026.
\bibitem{flashloop} FlashLoop contributors. FlashLoop: Fast and Memory-Efficient Looped Transformers via Lazy Updates. \href{https://github.com/Superone77/FlashLoop}{Official source}, frozen revision aeaadee.
\bibitem{megatron} NVIDIA. Megatron-LM and Megatron Core. \href{https://github.com/NVIDIA/Megatron-LM}{Official source}. Tested MCore revision 1dcf0da with the pinned VIME patch.
\bibitem{slime} Slime contributors. Slime: An LLM Post-Training Framework for RL Scaling. \href{https://github.com/THUDM/slime}{Official source}; formal looped baseline pending.
\bibitem{vime} VIME contributors. Native Looped-Model RL Integration. \href{https://github.com/vllm-project/vime/pull/464}{PR 464}; inherited source revision 71860e5.
\bibitem{rlt} vllm-rlt contributors. Native Training-Contract and Latent-Replay Implementation. \href{https://github.com/0z5a/vllm-rlt/tree/d2a358933393f25d74dd2bdd1068a741cd5d9226}{Inherited source d2a3589}.
\bibitem{artifact} 0z5a. ScaleRLT Implementation and Evidence Draft Stack. \href{https://github.com/0z5a/vime/pull/35}{Current preparation draft 35}, 2026. Individual source revisions and receipt digests are recorded below.
\bibitem{acceptance} 0z5a. Support vLLM-RLT as a Native Rollout Backend. \href{https://github.com/vllm-project/vime/issues/465}{VIME RFC 465}, 2026.
\bibitem{draftfeatures} 0z5a. Collect Versioned Target Features for External Draft. \href{https://github.com/vllm-project/vime/pull/454}{VIME PR 454}, merged October 6, 2026; source 76ad2be4, merge 5315305d.
\end{thebibliography}

\clearpage
\appendix
\section{Reproducibility Ledger}\label{app:ledger}
The evidence repository preserves immutable source/checkpoint identities, input and objective manifests, raw traces, all numerical failures, complete payload hashes and natural process completion/handback. Sources and historical devices are not merged into one execution record. Completed task-owned model files are removed only after their evidence is retained and no reader remains. Test counts from overlapping suites are not added.

\subsection{Checkpoint and source identities}
The ordinary inference checkpoint is ByteDance/Ouro-1.4B at \code{574fa66cb8bf5abdc979642d01cf2b79b16bfab1}; its 2,869,336,434-byte weight file has SHA-256
\begin{quote}\small\ttfamily\raggedright
58872a72616c736595b8b7662079c5b1\newline 2c5a162ec16eae94f21c348dfa9885af.
\end{quote}
The official Thinking input prepared for online RL is a different checkpoint: ByteDance/Ouro-1.4B-Thinking at \code{3aaa2224253a92ca45cf2e3d427c360e1ef9c93d}, weight SHA-256
\begin{quote}\small\ttfamily\raggedright
6730581fd366f6d1689728c616375202\newline 28385cdfc17c0b24f19aa840b00565fb.
\end{quote}
No ordinary-checkpoint inference result is promoted to Thinking online-learning evidence. IQuest is pinned at revision \code{61e8589747f6987ec7725e4ffe205f7a84561bd2} and has zero official tensor-body bytes downloaded in the reported header audit.

The official numerical/quality/cost baseline uses source e8aa1ee; historical three-family MCore uses 2ad271d; the repeated epoch2 B control uses d6229d6. Exact full revisions are retained in their manifests. MCore CUDA runs use the existing task-owned environment; no environment update is part of these experiments.

\begin{table*}[t]
\centering\scriptsize
\caption{Independently verified raw evidence anchors. Payload counts exclude the additional SHA manifest. Full digests identify immutable archives, not statistical replications.}
\label{tab:anchors}
\begin{tabular}{@{}p{0.25\textwidth}rp{0.61\textwidth}@{}}
\toprule
Archive & $n$ & SHA-256\\
\midrule
Three-family H20 MCore &84&\texttt{5cbea08375482e57936c9808a352fa3d12b22aabbb4a04a21723deac117f42ff}\\
Epoch2 B MCore &20&\texttt{d1dc4b22da2ce28d29e07106267ff5802134ab3faae4cade8d548dcec743e618}\\
Official C32 fixed-work &99&\texttt{06d960ebb83e97eb1dee8af3a98b5c3cc956815eb558c7cdabccc37077ec5cb4}\\
Matched compressed source &---&\texttt{772e1abe21ae85b6ccfa29ce03fc4e86cbe7c7082d3614a9fe397690fac2f043}\\
\bottomrule
\end{tabular}
\end{table*}

\subsection{Prepared data and accounting}
The public MATH reconstruction pins EleutherAI revision 21a5633 and MATH-500 revision 6e4ed1a. All 18 input files are hash verified. One normalized training duplicate and one overlap with the full test split are removed, leaving 7,498 training questions. All 500 held-out items remain unchanged. A fixed hash selects 256 development items, leaving 4,244 test questions reserved. Four explicit solution-hash-bound training label repairs and two unknown difficulty labels remain documented.

The source-prompt profile executes RLTT's unchanged conversion functions using actual pandas/PyArrow, preserving all 12,498 public item identities, labels and split membership. Two fresh preparations are byte-identical. The Thinking tokenizer's 1,024-token training cap excludes 19 explicitly listed questions in this profile, leaving 7,479; no held-out item is removed. Development/MATH-500 maximum prompt lengths are 1,013/972. This is a public-input reimplementation, not the authors' private dataset.

A separate integration repair prevents a rendered prompt from receiving a second chat template. All 500 official-tokenizer held-out prompts, two template modes and two completions retain prepared token IDs in 2,000 adapter checks. Known-answer transport doubles verify inputs and the grader, not generated reasoning quality.

Original RLTT batch semantics must be observed end to end. Conditional source inspection at CLI P32/G8 implies 256 prompt rows repeated eight times, or 2,048 new completions per rollout, rather than assuming 32 prompts. The declared independent-CLI profile instead uses 32 prompts and 256 completions. These are separate contracts; the unchanged full loader/FSDP/optimizer stack has not been executed. No conditional source count is presented as a measured training budget.

\subsection{Frozen startup and full-state requirement}
The prepared Thinking control uses seed 42, prompt cap 1,024, response cap 2,048, temperature 0.9, top-$p$ 1, natural termination, FP32 and response-mean RLTT. AdamW uses learning rate $10^{-6}$, betas 0.9/0.999, epsilon $10^{-8}$, weight decay 0.1, gradient clip 0.1 and KL coefficient 0.001. It keeps the ordinary MCore schedule with rematerialization and prefix sharing disabled. Three updates use four prompts and eight responses each; eight held-out completions at initial/three subsequent policies produce the 32-completion evaluation ledger.

Three full checkpoints and four physical publications require approximately 75 GiB of native state and a frozen minimum of 84 GiB private free storage. The earlier epoch2 input preflight verified 770 source files, all seven pinned Thinking-model files and 93.719 GiB free private XFS. A writable task-owned draft release and its published part quota are verified~\cite{artifact}. Neither this preflight nor a small source-archive upload establishes completed optimizer-state archival or recovery.

\section{Additional Component Results}\label{app:components}
\begin{table}[h]
\centering\small
\caption{Readout-only tiny CPU forward/backward. Saved bytes exclude supplied inputs and parameters. All speed ratios indicate regression.}
\begin{tabular}{@{}rrlrr@{}}
\toprule
Rows&Vocab&Objective&Ratio&Saved B\\
\midrule
64&1024&Score&0.414&$262144\to1024$\\
64&1024&Entropy/KL&0.444&$786432\to1024$\\
256&4096&Score&0.722&$4194304\to4096$\\
256&4096&Entropy/KL&0.654&$12582912\to4096$\\
\bottomrule
\end{tabular}
\end{table}

\begin{table}[h]
\centering\small
\caption{Tiny CPU rematerialization versus untouched parent. All 95\% intervals include one. Saved storage is a separate forward-hook metric.}
\begin{tabular}{@{}lrrr@{}}
\toprule
Family&Ratio&95\% interval&Saved B\\
\midrule
Ouro&1.119&[0.514,2.436]&$51360\to8512$\\
Nanbeige&0.649&[0.268,1.570]&$27688\to6008$\\
Huginn&0.517&[0.261,1.022]&$71660\to12252$\\
\bottomrule
\end{tabular}
\end{table}

Forward saved storage counts unique non-parameter/non-buffer backing addresses, including complete view storage, and excludes backward workspace and allocator overhead. A separate joint replay observation reduces prompt saved storage from 24,544 to 32 bytes while mandatory boundary storage remains 2,560 bytes and decoder token calls double from 80 to 160. These counts are not official-model memory or speed measurements.

The cold G8/C8 rollout fork achieves tiny CPU ratios 1.811/3.631/3.622 for Ouro P17/128/129, and 1.828/4.510/2.570 for Nanbeige. The last interval is [0.259,25.497]. One-request controls remain slower, and snapshot retention can increase live pages. All 3,504 measured requests match serial oracles in tokens, depth, stops and versions. Both arms allocate the same page pool. These results do not establish lower allocated GPU memory.

\subsection{Evidence transport}
Three source-archive routes use the exact same 10,178,626 compressed bytes, including upload, independent download, server-digest verification and control requests (Table~\ref{tab:transport}). Each route has only one timed attempt, so no replication-based interval is reported. These descriptive ratios do not measure model download, GPU kernels or RL.

\begin{table}[h]
\centering\small
\caption{Matched small source-evidence transport. Ratio is relative to direct TLS; the archived source digest is in Table~\ref{tab:anchors}.}
\label{tab:transport}
\begin{tabular}{@{}lrr@{}}
\toprule
Route&Seconds&Ratio vs direct\\
\midrule
TLS via H20 SOCKS&68.670&$0.352\times$\\
Direct local TLS&24.183&$1.000\times$\\
Existing Mac HTTP proxy&6.761&$3.577\times$\\
\bottomrule
\end{tabular}
\end{table}

The separate uncompressed 38,015,488-byte source transfer takes 84.780 seconds and is not a matched input for these ratios. Its historical 75 GiB serial extrapolation is not a current-route measurement. The current small-upload/writeability and part-quota checks pass; an actual 75 GiB full-state transfer remains unrun.

\section{Claim-to-Evidence Map}
The accompanying Markdown result table maps each claim to its original report and immutable receipt. The primary empirical results use the independently audited archives and full-gradient references listed here. Short runs under other objectives do not establish this system's reward convergence. New measurements replace pending rows only after their own model, worker, correctness, recovery and complete-state gates pass.
\end{document}
